\section{The internal advantage by kind of error}
\label{sec:result}

\begin{table*}[t]
\centering
\footnotesize
\setlength{\tabcolsep}{1.5pt}
\caption{Every run under the corrected labels. $n_+$ and $n_-$ are the labelled failures and valid calls in the scored population. The AUC columns are the length floor, the mean log-probability (conf.), the output-only judge, the Qwen3.5-0.8B reader and the token-role probe, all on every scored failure. The last four columns are differences of AUC on held-out tools: probe minus confidence on every scored failure with its 95\% tool-resampled interval, probe minus the output-only judge on every scored failure, and probe minus confidence on wrong argument values and on dropped parallel calls within the parallel categories, each against valid calls. $^{\circ}$ the interval excludes zero (intervals in \cref{app:judges,app:types}). n.c.\ not computed: too few failures of that kind or of valid parallel calls, an underpowered run, or a comparator not fitted on that run. $^{*}$ re-extracted under the corrected extractor. $^\dagger$ fewer than thirty items of one class.}
\label{tab:main}
\begin{tabular}{@{}llrrccccccccc@{}}
\toprule
Model & Data & $n_+$ & $n_-$ & floor & conf. & judge & reader & probe & probe $-$ conf. & probe $-$ judge & value err. & dropped \\
\midrule
Llama-3.2-1B & Glaive & 221 & 359 & 0.539 & 0.576 & 0.910 & 0.949 & 0.963 & $+0.387$ {\scriptsize$[+0.18, +0.59]$} & $+0.053^{\phantom{\circ}}$ & $+0.357^{\circ}$ & n.c. \\
Llama-3.2-1B$^{*}$ & BFCL & 161 & 339 & 0.635 & 0.651 & 0.755 & 0.818 & 0.831 & $+0.180$ {\scriptsize$[+0.09, +0.27]$} & $+0.077^{\circ}$ & $+0.240^{\circ}$ & n.c. \\
Llama-3.2-1B & BFCL-live & 208 & 157 & 0.709 & 0.663 & 0.682 & 0.795 & 0.782 & $+0.119$ {\scriptsize$[-0.00, +0.24]$} & $+0.101^{\circ}$ & $-0.010^{\phantom{\circ}}$ & n.c. \\
Llama-3.2-3B & Glaive & 90 & 555 & 0.522 & 0.785 & 0.684 & 0.931 & 0.911 & $+0.126$ {\scriptsize$[+0.02, +0.28]$} & $+0.227^{\phantom{\circ}}$ & $+0.145^{\circ}$ & n.c. \\
Llama-3.2-3B & BFCL & 120 & 576 & 0.618 & 0.707 & 0.633 & 0.758 & 0.831 & $+0.125$ {\scriptsize$[+0.05, +0.20]$} & $+0.199^{\circ}$ & $+0.134^{\circ}$ & $+0.158^{\circ}$ \\
Gemma-3-1B & Glaive & 158 & 460 & 0.704 & 0.650 & 0.872 & 0.853 & 0.883 & $+0.233$ {\scriptsize$[+0.11, +0.34]$} & $+0.011^{\phantom{\circ}}$ & $+0.226^{\circ}$ & n.c. \\
Gemma-3-1B$^{*}$ & BFCL & 364 & 150 & 0.842 & 0.711 & 0.900 & 0.927 & 0.948 & $+0.237$ {\scriptsize$[+0.19, +0.29]$} & $+0.048^{\circ}$ & $+0.158^{\phantom{\circ}}$ & $+0.342^{\circ}$ \\
\midrule
Qwen3-1.7B & BFCL & 52 & 639 & 0.628 & 0.799 & 0.558 & 0.662 & 0.745 & $-0.053$ {\scriptsize$[-0.16, +0.05]$} & $+0.187^{\circ}$ & $-0.124^{\phantom{\circ}}$ & n.c. \\
Qwen3-1.7B$^\dagger$$^{*}$ & Glaive & 10 & 699 & 0.677 & 0.867 & n.c. & n.c. & 0.676 & $-0.191$ {\scriptsize$[-0.65, +0.21]$} & n.c. & n.c. & n.c. \\
MiniCPM5-2B$^{*}$ & BFCL & 54 & 640 & 0.578 & 0.823 & 0.585 & 0.677 & 0.775 & $-0.048$ {\scriptsize$[-0.15, +0.05]$} & $+0.190^{\circ}$ & $-0.162^{\circ}$ & n.c. \\
MiniCPM5-2B$^{*}$ & BFCL-live & 120 & 475 & 0.545 & 0.797 & 0.642 & 0.655 & 0.744 & $-0.053$ {\scriptsize$[-0.15, +0.02]$} & $+0.102^{\circ}$ & $-0.048^{\phantom{\circ}}$ & n.c. \\
MiniCPM5-2B$^\dagger$$^{*}$ & Glaive & 9 & 716 & 0.685 & 0.913 & n.c. & n.c. & 0.710 & $-0.203$ {\scriptsize$[-0.51, +0.01]$} & n.c. & n.c. & n.c. \\
Qwen3.5-0.8B$^{*}$ & BFCL & 86 & 398 & 0.601 & 0.792 & 0.678 & 0.730 & 0.682 & $-0.110$ {\scriptsize$[-0.19, -0.03]$} & $+0.004^{\phantom{\circ}}$ & $-0.121^{\circ}$ & n.c. \\
\midrule
MiniCPM5-2B$^{*}$ & BFCL, JSON & 83 & 514 & n.c. & 0.757 & n.c. & n.c. & 0.882 & $+0.125$ {\scriptsize$[+0.05, +0.19]$} & n.c. & $-0.066^{\phantom{\circ}}$ & $+0.201^{\circ}$ \\
Qwen3.5-0.8B$^{*}$ & BFCL, JSON & 226 & 436 & n.c. & 0.808 & n.c. & n.c. & 0.868 & $+0.060$ {\scriptsize$[+0.01, +0.11]$} & n.c. & $-0.136^{\circ}$ & $+0.171^{\circ}$ \\
\bottomrule
\end{tabular}

\end{table*}

On every scored failure, the \nRunsInternals{} powered Llama-3.2 and Gemma-3 runs have an internal advantage of \gapInternalsMin{} to \gapInternalsMax{}, with an interval excluding zero on \nInternalsCiExcludesZero{} of them (\cref{tab:main}, \cref{fig:types}a). The \nRunsConfidence{} powered Qwen3, Qwen3.5 and MiniCPM5 runs read \gapConfidenceMin{} to \gapConfidenceMax{}, with an interval below zero on \nConfidenceCiBelowZero{}. Runs that share a checkpoint are not independent. With \nCkptInternals{} against \nCkptConfidence{} checkpoints the smallest attainable two-sided rank-test value is $p=\ckptExactFloor{}$, so the separation is a description of these \nCkpt{} checkpoints.

\paragraph{Wrong argument values.}
Every scored failure mixes kinds of error, so \cref{fig:types}b scores wrong argument values against valid calls alone. The probe leads confidence on \nWavInternalsPos{} of the \nInternalsRuns{} probe-side runs, by \wavInternalsMin{} to \wavInternalsMax{} in point estimate, with an interval excluding zero on \nWavInternalsCiExcl{}. The largest lead, Llama-3.2-1B on Glaive, has the widest interval, [\WavGapLoLlamaOneBGlaive{}, \WavGapHiLlamaOneBGlaive{}]. Llama-3.2-1B on the live categories reads \WavGapLlamaOneBLive{}, on the probe side by family and on neither side by measurement. Confidence leads the probe on all \nWavConfidenceRuns{} confidence-side runs, native and forced into JSON, where the advantage is \wavConfidenceMin{} to \wavConfidenceMax{}, with an interval excluding zero on \nWavConfidenceCiBelow{}. On these models the mean log-probability ranks most wrong values below right ones. Glaive repeats requests, and keeping the first occurrence of each unique request leaves the value-error advantage at \gdWavLlamaOneBGlaive{} on Llama-3.2-1B and \gdWavGemmaGlaive{} on Gemma-3.

\paragraph{Dropped parallel calls.}
A one-bit flag for a request in a parallel category separates dropped calls from valid calls at \catAucMin{} to \catAucMax{} AUC on the four runs with at least twenty of them (\cref{fig:types}c), because only parallel requests can drop a call. On Gemma-3 the flag reaches \catAucGemmaBfcl{} against the probe's \catProbeGemmaBfcl{}. Three of these runs use the fallback prompt that asks for one JSON object. Within the parallel categories, where the flag is constant, the probe leads confidence on MiniCPM5-2B and Qwen3.5-0.8B under forced JSON, \PwGapMiniCpmBfclJson{} [\PwGapLoMiniCpmBfclJson{}, \PwGapHiMiniCpmBfclJson{}] and \PwGapQwenThreeFiveBfclJson{} [\PwGapLoQwenThreeFiveBfclJson{}, \PwGapHiQwenThreeFiveBfclJson{}]. On Llama-3.2-3B, the only one of the four prompted through its native template, the lead vanishes at \PwGapLlamaThreeBBfcl{} [\PwGapLoLlamaThreeBBfcl{}, \PwGapHiLlamaThreeBBfcl{}], and Gemma-3 makes too few valid parallel calls to be tested. No claim about dropped calls is made beyond these three numbers.

\paragraph{Output format and R3.}
Qwen3 writes JSON like Llama-3.2 and Gemma-3 and sits on the confidence side, so the native format does not decide the side. Registered prediction R3 said that, forced to write JSON, the confidence-side families keep an advantage at or below zero on every scored failure, and that it fails above \rThreeThreshold{}. It failed on MiniCPM5-2B, at \rThreeJsonMiniCpm{} [\rThreeJsonLoMiniCpm{}, \rThreeJsonHiMiniCpm{}], on a partial run of \nJsonItemsMiniCpm{} of \nAllMiniCpmBfcl{} items extracted from a tree with uncommitted changes. On Qwen3.5-0.8B it is undecided by the rule, at \rThreeJsonQwenThreeFive{} [\rThreeJsonLoQwenThreeFive{}, \rThreeJsonHiQwenThreeFive{}]. Dropped parallel calls make up \rThreeMcJsonMiniCpm{} of \rThreePosJsonMiniCpm{} and \rThreeMcJsonQwenThreeFive{} of \rThreePosJsonQwenThreeFive{} failures under the fallback prompt, against \rThreeMcNatMiniCpm{} and \rThreeMcNatQwenThreeFive{} under the native templates. On wrong argument values both models stay on the confidence side.

\paragraph{Size and release date.}
Size does not separate the sides (one and three billion parameters against 0.8 to two), and within Llama-3.2 the probe does not improve from 1B to 3B (\aucProbeLlamaOneBBfcl{} at both sizes on BFCL). Release date, a reasoning mode and native tool-calling post-training all separate the six checkpoints (\cref{sec:discussion}).

\paragraph{Calls with a history.}
One multi-turn run, MiniCPM5-2B, reads \mtAucProbe{} for the probe and \mtAucConf{} for confidence, the opposite order to its single-turn runs, on \mtGroupsAll{} linked groups. It is anecdotal and enters no claim (\cref{app:multiturn}).
